\documentclass[11pt]{article}
\usepackage[a4paper, margin=2.5cm]{geometry}
\usepackage{amsmath, amssymb, amsthm}
\usepackage{booktabs}
\usepackage{hyperref}
\usepackage{graphicx}
\usepackage{xcolor}

\title{On the $\sqrt{R}$ measurement-error misspecification\\
in BGAM for trial-aggregated EEG data}
\author{Jiahui Fan}
\date{May 22, 2026}

\begin{document}
\maketitle

\begin{abstract}
The meta-analytic formulation of Nalborczyk \& B\"urkner (2025)
fits a Bayesian GAM to trial-aggregated EEG data via
\texttt{eeg\_mean | se(eeg\_sd) \textasciitilde\, ...}.
We show that this specification systematically inflates the
measurement-error standard deviation by a factor of $\sqrt{R}$,
where $R$ is the number of trials per cell. The bias is mild in
the original 1D smoothing setting but becomes dominant in the
3D tensor-product extension developed in this thesis, contributing
substantially to the observed peak-amplitude under-recovery.
We derive the correct specification, quantify the bias, and propose
a one-line fix.
\end{abstract}

\section{The meta-analytic specification}

The raw EEG simulation has dimensions
\begin{align*}
\text{trials}\,(R) \times \text{participants}\,(P)
\times \text{channels}\,(C) \times \text{time}\,(T)
\times \text{conditions}\,(K) ,
\end{align*}
giving $RPCTK$ rows in total. Fitting BGAM on raw rows is
computationally infeasible (millions of observations $\times$
thousands of tensor basis coefficients). The standard
preprocessing step (\texttt{aggregate\_trials()} in our utility
module) aggregates each $(p, c, k, t)$ cell across its $R$ trials:
\begin{align*}
\bar y_{pcktk} &= \frac{1}{R}\sum_{r=1}^{R} y_{rpck t},\\
s_{pck t}     &= \sqrt{\frac{1}{R-1}\sum_{r=1}^{R}
                 (y_{rpck t} - \bar y_{pck t})^2}.
\end{align*}
Following Nalborczyk \& B\"urkner (2025), the model is then fit
to the cell-level summaries with measurement error declared
through the \texttt{brms} \texttt{se()} term:
\begin{equation}
\texttt{eeg\_mean | se(eeg\_sd) \textasciitilde\,
condition + (1 + condition | participant) + smooth term}.
\label{eq:nb-formula}
\end{equation}

\section{What \texttt{brms se()} actually does}

The \texttt{se()} argument tells \texttt{brms} that the outcome
variable carries known measurement noise with the supplied
standard deviation. Concretely, the likelihood becomes
\begin{equation}
y_{\text{obs}} \,\sim\, \mathcal{N}\!\bigl(\mu,\,
\sigma_{\text{known}}^{\,2} + \sigma_{\text{res}}^{\,2}\bigr),
\label{eq:brms-likelihood}
\end{equation}
where $\sigma_{\text{known}}$ is the value passed into
\texttt{se(.)} and $\sigma_{\text{res}}$ is an additional residual
SD estimated by \texttt{brms}. The interpretation of
$\sigma_{\text{known}}$ is therefore the \emph{standard deviation
of the outcome variable as a random variable}, not the dispersion
of any auxiliary quantity.

\section{The discrepancy}

For our setting, $y_{\text{obs}} = \bar y_{pck t}$ is the
\emph{mean} of $R$ trials, not an individual trial. The
distribution of $\bar y$ around its expectation is governed by
the central limit theorem:
\begin{equation}
\operatorname{Var}(\bar y_{pck t}) \,=\, \frac{\sigma_{\text{trial}}^{\,2}}{R},
\qquad
\operatorname{SD}(\bar y_{pck t}) \,=\, \frac{\sigma_{\text{trial}}}{\sqrt{R}}.
\label{eq:clt}
\end{equation}
The natural sample-based estimator of this
\emph{standard error of the mean} is
\begin{equation}
\widehat{\operatorname{SE}}(\bar y_{pck t})
   \,=\, \frac{s_{pck t}}{\sqrt{R}}
   \,=\, \texttt{eeg\_sem}.
\end{equation}
However, the formulation \eqref{eq:nb-formula} passes
$s_{pck t} = \texttt{eeg\_sd}$ rather than
$s_{pck t}/\sqrt{R} = \texttt{eeg\_sem}$. By
\eqref{eq:brms-likelihood}, this inflates the measurement-error
SD that \texttt{brms} uses for $\bar y$ by a factor of $\sqrt{R}$:
\begin{equation}
\boxed{\;
\frac{\text{value passed to \texttt{se()}}}
     {\text{correct SE of the mean}}
\;=\;
\frac{s_{pck t}}{s_{pck t}/\sqrt{R}}
\;=\;\sqrt{R}.\;}
\end{equation}

For our pilot data with $R=20$, this is a factor of
$\sqrt{20}\approx 4.47$.

\section{Effect on the posterior smooth}

The fitted smooth surface $\hat f$ minimises (up to the prior)
the negative log-likelihood
\begin{equation}
\ell(f) \,=\, \frac{1}{2}\sum_{\text{cells}}
\frac{(\bar y - f)^2}{\sigma_{\text{known}}^{\,2} + \sigma_{\text{res}}^{\,2}}
+ \text{smoothness penalty}.
\end{equation}
Inflating $\sigma_{\text{known}}$ by $\sqrt{R}$ inflates the
denominator by a factor of up to $R$, which reduces the
\emph{relative} weight of every data residual against the
smoothness penalty by the same factor. The posterior therefore
shifts toward the smoothness prior --- a smoother, lower-amplitude
surface --- regardless of the actual signal magnitude. This is
the mechanism by which the misspecified \texttt{se()} produces
systematic under-recovery of the peak amplitude.

\section{Why 1D masks the bias and 3D exposes it}

The strength of the smoothness prior relative to the likelihood
scales with the dimensionality of the basis. Let $d$ denote the
total number of tensor basis coefficients constrained by the data.
For $N$ effective data points, the prior-to-likelihood ratio is
approximately $d/N$.

\begin{itemize}
\item[\textbf{1D (N\&B 2025):}]
$d = k_t = 20$. Even with the $\sqrt{R}$ inflation, the
likelihood remains the dominant term, and the resulting peak
under-recovery is on the order of 10--20\%.

\item[\textbf{3D (this thesis):}]
$d = k_{xy} \cdot k_t \cdot K \cdot 2 = 20 \cdot 20 \cdot 2 \cdot 2
\approx 1600$ coefficients (using full ANOVA decomposition).
With only $P=5$ participants and $R=20$ trials per cell, the
likelihood becomes weak relative to the prior, and the
$\sqrt{R}$-inflated noise tilts the posterior decisively toward
over-smoothing. Empirically, peak under-recovery reaches 54\%.
\end{itemize}

The same misspecification therefore has qualitatively different
consequences in the two settings: a mild bias in 1D, a dominant
bias in 3D.

\section{Empirical verification}

Table~\ref{tab:recovery} reports peak-amplitude recovery across
fits that vary in the SE specification while holding all other
hyperparameters fixed.

\begin{table}[h]
\centering
\small
\begin{tabular}{lllcc}
\toprule
Fit                  & Algorithm & SE term         & Peak ($\mu$V) & Recovery (\%) \\
\midrule
v5\_HMC\_baseline    & HMC       & \texttt{se(eeg\_sd)}  & $0.69$ & 46 \\
A\_no\_ME (Exp.\ A)  & VI        & none                  & $1.18$ & 79 \\
v5\_baseline\_VI     & VI        & \texttt{se(eeg\_sd)}  & $\cdot$ & $\cdot$ \emph{(in-flight)} \\
v5\_k3030\_sem       & HMC       & \texttt{se(eeg\_sem)} & $\cdot$ & $\cdot$ \emph{(in-flight)} \\
v5\_k3030\_nose      & HMC       & none                  & $\cdot$ & $\cdot$ \emph{(in-flight)} \\
\bottomrule
\end{tabular}
\caption{Peak-amplitude recovery under different SE specifications.
Removing the SE term entirely (Exp.\ A) recovers 79\% of the
truth, compared with 46\% under the original
\texttt{se(eeg\_sd)} specification. The \texttt{se(eeg\_sem)}
variant, currently in-flight, is predicted to recover 70--85\%
based on the $\sqrt{R}$ analysis above.}
\label{tab:recovery}
\end{table}

\section{Proposed fix}

The fix is a single-token change. In
\texttt{code/utils.R}, the aggregator now emits both
\texttt{eeg\_sd} and \texttt{eeg\_sem}\,$=\,$\texttt{eeg\_sd/sqrt(R)};
in \texttt{code/01\_fit\_bgam.R}, the \texttt{se\_mode} argument
selects between three formulations:
\begin{itemize}
\item \texttt{se\_mode = "sd"}: original N\&B recipe (kept for
  back-compatibility).
\item \texttt{se\_mode = "sem"}: $\sqrt{R}$-corrected SE
  (\textbf{recommended}).
\item \texttt{se\_mode = "none"}: no measurement-error term
  (matches Exp.\ A).
\end{itemize}
The three variants will be benchmarked on identical data; the
\texttt{se(eeg\_sem)} fit is the principal correction.

\section{Discussion and relation to prior work}

The misspecification analysed here arises from conflating the
\emph{trial-level variability} ($\sigma_{\text{trial}}$, estimated
by \texttt{eeg\_sd}) with the \emph{measurement-error SD of the
mean estimate} ($\sigma_{\text{trial}}/\sqrt{R}$, estimated by
\texttt{eeg\_sem}). The two are conceptually adjacent --- both
describe ``how noisy is the data?'' --- but quantitatively
differ by $\sqrt{R}$ for $R$-trial averages. Standard
meta-analytic practice (Hedges \& Olkin 1985; Borenstein et al.\
2009) consistently passes the standard error of the effect
estimate, not the within-study SD.

In the original \texttt{brms} documentation (B\"urkner 2017), the
\texttt{se()} argument is defined as the ``known measurement error
standard deviation'', which by \eqref{eq:brms-likelihood}
unambiguously refers to the SD of the outcome. The
\texttt{se(eeg\_sd)} usage thus represents a subtle conceptual
slip that is easy to overlook in 1D applications, where its
practical consequences are small.

Our 3D extension surfaces this slip empirically. The proposed
$\sqrt{R}$ correction does not change the spirit of the N\&B
meta-analytic framework; it adopts the standard meta-analytic
convention for the SE of the mean, restoring agreement with the
implicit \texttt{brms} likelihood.

\section{Conclusion}

We have identified that the BGAM specification
\texttt{eeg\_mean | se(eeg\_sd)} of Nalborczyk \& B\"urkner (2025)
inflates the meta-analytic standard error of the cell mean by a
factor of $\sqrt{R}$. The bias is benign in the 1D smoothing
setting of the original paper but dominates in the 3D tensor
extension developed here, accounting for a substantial fraction
of the observed peak-amplitude under-recovery. The proposed fix,
\texttt{se(eeg\_sem)} with
\texttt{eeg\_sem = eeg\_sd / sqrt(R)}, is a one-line change with
no other modeling implications and recovers consistency with the
standard \texttt{brms} likelihood.

\section*{References}

\begin{itemize}
\item Borenstein, M., Hedges, L.~V., Higgins, J.~P.~T., \&
  Rothstein, H.~R. (2009). \emph{Introduction to meta-analysis}.
  Wiley.
\item B\"urkner, P.-C. (2017). brms: An R package for Bayesian
  multilevel models using Stan. \emph{Journal of Statistical
  Software}, 80(1).
\item Hedges, L.~V., \& Olkin, I. (1985). \emph{Statistical
  methods for meta-analysis}. Academic Press.
\item Nalborczyk, L., \& B\"urkner, P.-C. (2025).
  Bayesian Generalized Additive Models for EEG. \emph{[venue]}.
\end{itemize}

\end{document}
